\section{Introduction}
Modern electronic design automation (EDA) flows coordinate synthesis, floorplanning, placement, clock-tree synthesis, routing, extraction, and physical verification. Building learning datasets from public RTL adds a second layer of work: recovering compilation closure, resolving implementation failures, aligning circuit entities across stages, and extracting physical labels with explicit prediction-time availability. Dataset production therefore depends on a sequence of decisions about which sources to accept, how to recover failed runs, and when the resulting graphs satisfy the intended learning contract.

OpenROAD provides physical implementation infrastructure~\citep{openroad}; ForgeEDA and HighTide support circuit collection and evolving RTL-to-GDS benchmarks~\citep{forgeeda,hightide}; AiEDA and CircuitOps connect physical-design artifacts to programmatic data access~\citep{aieda,circuitops}. Building on these capabilities, we address the evidence handoffs between acquisition, physical recovery, and graph admission. Each handoff has a distinct acceptance condition: reproducible synthesis for a source candidate, validated outcomes for a reusable repair, and consistent identities, topology, and labels for a graph sample.

We present \emph{Agentic RTL-to-Graph}, implemented by the R2G skills workflow. Structured run observations guide candidate selection, scoped repair, stage resumption, and admission; deterministic programs execute these actions and preserve their provenance. A candidate-to-recipe lifecycle turns development-time repair proposals into policies that can be reused without further LLM calls. A four-stage graph contract anchors entities to a canonical synthesis netlist and separates permitted input features from post-route supervision. Independent semantic checks make these requirements measurable at the output boundary. This matters for learning because a correct connectivity graph can still lack the physical features and target values needed to define a supervised sample. R2G treats entity identity, physical quantities, stage visibility, and validity masks as one conversion contract.

Our contributions are fourfold. First, we organize RTL-to-graph construction around explicit source, run, policy, and graph evidence contracts. Second, we implement and evaluate frozen repair reuse with disjoint proposal, validation, and held-out designs. Third, we construct four-stage graphs with independent semantic acceptance checks that cover identity, topology, numerical features, labels, masks, and causal stage restrictions. Fourth, we show on a fixed 278-design split that the admitted graphs retain useful stage-specific physical signal: placement geometry sharply improves wirelength and effective-resistance prediction, while graph topology helps selected targets rather than universally dominating feature-only models. Complementary experiments also demonstrate 200 qualified RTL designs from 122 repositories, eight repairs on 13 held-out challenges using the candidate bank, and fully accepted graph outputs for all 24 test conversion designs. Throughout the paper, ``R2G'' denotes our recorded implementation label; ``R2G benchmark'' denotes the prior multi-view dataset~\citep{r2gbenchmark}.

\subsection{Relation to Prior Work}
\textbf{Datasets and representations.} CircuitNet and OpenABC-D provide data for physical-design and synthesis learning~\citep{circuitnet,openabcd}. ForgeEDA supplies RTL, mapped netlists, AIGs, and placed netlists~\citep{forgeeda}. AiEDA and CircuitOps provide flow-to-data and graph-oriented infrastructure~\citep{aieda,circuitops}; EDA-Schema-V2 organizes stage-resolved physical data across technologies~\citep{edaschemav2}. The R2G benchmark standardizes five graph views for representation comparisons~\citep{r2gbenchmark}. Our work focuses on qualified sample production and evaluates a four-stage prediction product with explicit identity, visibility, and label checks.

\textbf{Flow agents and evolving benchmarks.} ChatEDA and ORFS-agent study tool orchestration and parameter optimization~\citep{chateda,orfsagent}. HighTide uses agent skills, incremental RTL-to-GDS compilation, and persistent per-design rationale for evolving curation~\citep{hightide}. Tool-interactive RTL-to-GDS benchmarks further evaluate complete EDA execution~\citep{rtltogdsbench}. Our emphasis is the connection between validated repair adoption and run-consistent graph admission, with separate measurements of acquisition yield, held-out repair outcomes, and semantic output quality.

\textbf{Repair evaluation.} PostEDA-Bench studies DRC fixing and PPA convergence with machine-checkable outcomes~\citep{postedabench}. Our repair experiment examines a complementary setting: proposals are learned and validated before deployment, then reused on held-out designs under a fixed action policy. Separating these phases measures the transfer of executable repair knowledge and its deployment cost.
